% titel: Blatt 9, A5: Begriffe definieren
% Teil II: die Definitionen.

\subsection*{a) K$_c$ und K$_p$}
\textbf{K$_c$} ist die Gleichgewichtskonstante, die mit molaren
Konzentrationen gebildet wird:
$K_c = [C]^c[D]^d/([A]^a[B]^b)$.
\textbf{K$_p$} ist dieselbe Konstante, aber mit Partialdrücken der Gase:
$K_p = p_C^c p_D^d / (p_A^a p_B^b)$. Der Zusammenhang $K_p = K_c(RT)^{\Delta n}$
($\Delta n = (c+d)-(a+b)$, Summe der Gaskoeffizienten) verknüpft beide.

\subsection*{b) Prinzip des kleinsten Zwanges}
Wird ein chemisches Gleichgewicht einem Zwang ausgesetzt (Änderung von
Druck, Temperatur oder Konzentration), verschiebt es sich so, dass der
äußere Zwang vermindert wird. Beispielsweise: mehr Edukt-Einfluss
verschiebt zur Produktseite; Temperaturerhöhung begünstigt die
endotherme Richtung.

\subsection*{c) Dynamisches Gleichgewicht}
Der Zustand, in dem Hin- und Rückreaktion gleich schnell ablaufen
($v_{\text{hin}}=v_{\text{rück}}$), sodass sich die Konzentrationen nicht
mehr ändern --- obwohl die Reaktionen auf Teilchenebene ständig weiterlaufen.

\subsection*{d) Reaktionsquotient Q}
Der Ausdruck der Massenwirkungsgleichung, aber mit den \emph{aktuellen}
Konzentrationen zu einem beliebigen Zeitpunkt.
Vergleich mit $K_c$: $Q<K$ (Gemisch reagiert zu Produkten), $Q=K$
(Gleichgewicht), $Q>K$ (reagiert zu Edukten).

\begin{ergebnis}
$K_c$/K$_p$ = Gleichgewichtskonstante (Konz. bzw. Partialdrücke);
Le Chatelier-Prinzip = Reaktion gegen den Zwang; dynamisches
Gleichgewicht = gleiche Vor- und Rückgeschwindigkeit; Q =
Reaktionsquotient, Richtungsanzeiger.
\end{ergebnis}

\begin{kontrolle}
Alle vier Definitionen gehören zusammen: nur $T$ verändert $K$; die
Lage steuert man mit $c$, $p$ oder $T$, die Richtung liest man an
$Q$ vs.\ $K$.
\end{kontrolle}

\begin{merksatz}
K ist eine Konstante bei fester $T$. Q sind die gleichen Quotienten zu
nicht-Gleichgewichtskonzentrationen. Le Chatelier erklärt jede
Verschiebung, dynamisches Gleichgewicht den stationären Zustand.
\end{merksatz}
